\chapter{Simulation Study}
In this section, we present a simulation study designed to evaluate the performance 
of the two methods for estimating causal effects with time-varying treatments 
under time-invariant unmeasured confounding. The data generation process (DGP) is taken 
from the DGP described in Section 5.1 of \cite{blackwellEffectPoliticalAdvertising2024}. In this 
DGP, the scenario which we simulate is a study design where individuals are exposed to a binary treatment
(e.g., receiving a political advertisement) over a period fixed period of time, and the outcome of interest
is measured only once at the end of the study period (e.g., voting behavior). The treatment assignment
is influenced by both observed covariates and an unmeasured time-invariant confounder (e.g., political
leaning). The outcome is affected by the treatment received at the final time point, the cumulative 
effect of treatments received in the last three time points, the mean of observed covariates over time, 
and the unmeasured confounder. 

We simulate longitudinal data for $n \in \{200, 500, 1000, 3000\}$ individuals. For the
length of the time serieses for each n, we consider three different lenghts: $\rho = n/T$ 
with $\rho \in \{5,10,50\}$. The data is balanced. menaning that each individual has 
the same number of time points. In Eq.~\eqref{eq:blackwell_scm}, we present 
the structural causal equations for the DGP. Here, $\operatorname{expit}(x) = \frac{1}{1 + e^{-x}}$ and 
$\bar{X}_{iT} = \frac{1}{T}\sum_{t=1}^T X_{it}$.

\begin{subequations}\label{eq:blackwell_scm}
\begin{align}
U_i &:= a(2V_i-1), \qquad V_i \sim \operatorname{Unif}(0,1)
\\
X_{it} &:= -\tfrac12\mathbbm{1} + L \xi_{it} \qquad \xi_{it}\sim \mathcal{N}(0,I), \qquad LL^\top=\Sigma
\\
A_{it} &:= \mathbbm{1}\!\left\{\eta_{it} \le
\operatorname{expit}\!\left(U_i+\varphi A_{i,t-1}+\beta^\top X_{it}\right)\right\},
\qquad \eta_{it} \sim \operatorname{Unif}(0,1)
\\
Y_{iT} &:= U_i + \tau_F A_{iT} + \tau_C \sum_{t=T-3}^{T-1} A_{it}
+ \gamma^\top \bar X_{iT} + \epsilon_{iT},
\qquad \epsilon_{iT} \sim \mathcal{N}(0,1)
\end{align}
\end{subequations}

We simulate for two different values of $a \in \{1, 2\}$ \hl{to control the
contribution of the unmeasured confounder $U_i$ to the variance of the linear predictor.}. For the number 
of covariates, we simulate for $p \in \{2, 4\}$. 
The coefficients are set to $\beta = (-0.5, -0.5)^{\top}$ when $p=2$ and 
$\beta = (-0.5, -0.5, 1, -0.5)^{\top}$ when $p=4$. \hl{(add some note about why we chose 
these coeffs based on the variance or sth)}. We set $\varphi = 0.3$. Note that we generate 
the covariates $X_{it}$ exogenously to the treatments. So we do not have the 
treatment-confounder feedback. This is done to isolate the effect of
adjusting for unmeasred time-invariant confounding. Also note that theere is no dynamics 
in the covariates and the distribution of $X_{it}$ is the same for all time points. The parameters 
for the outcome model are set to $\tau_F = 1$, $\tau_C = 0.3$, and $\gamma = (1, 0.5)^{\top}$ 
when $p=2$ and when $p=4$, $\gamma = (1, 0.5, 1, 1)^{\top}$.

Our target estimands are the average treatment of the last treatment: 
$\mathbb{E}\!\left[ Y_i\!\left({A}_{T} = 1\right)  -  Y_i\!\left({A}_{T} = 0\right) \right]$ and the average 
treatment effect of changing one of the last three treatments from 0 to 1: 
$\mathbb{E}\!\left[ Y_i\!\left(\underline{d}_{T-k}\right) \right]$,
 and  and the cumulative effect of the last three treatments:



